% Solution: Mas-Colell, Whinston & Green (1995) Exercise 8.B.5. Input by \solitem.
% Reused from lista-reais-lectures-01-04-solutions (MWG 8.B.5), deepened. Check: check_A1.py.
\begin{exsol}{MWG Exercise 8.B.5}{262}{Advanced}{5}
\asked Show that IESDS gives a unique prediction in linear Cournot with two firms, and say whether
it still does with three.

\method Iterated deletion of strictly dominated strategies (MWG \S8.B, Defs.~8.B.2--8.B.4 and the
iterated procedure, pp.~236--241; Lecture~1, slide~29), applied to a continuum of quantities by
bounding best replies, as in Tadelis \S4.2.2, pp.~65--67, and the slides' Cournot remark
(Lecture~2, slide~24: ``eliminating outputs that are never optimal yields the sequence
$\tfrac12 q^m, \tfrac34 q^m, \tfrac58 q^m, \dots \to \tfrac23 q^m$''). The tool is the dominance
lemma of the roommates exercise (Tadelis 4.6), here with a slope $b$. Assumptions: $a > c \ge 0$,
$b > 0$, $q_i \in [0, \infty)$; write $m = (a-c)/(2b)$ for the monopoly output.

\solution
\textbf{Payoff and best reply.} With $Q_{-i}$ the rivals' total output,
\[
  u_i = \big(a - b(q_i + Q_{-i}) - c\big)\, q_i
      = b\,q_i\Big(\frac{a-c}{b} - Q_{-i} - q_i\Big)
      = b\, q_i\,(2x - q_i), \qquad x = m - \tfrac{Q_{-i}}{2},
\]
since $(a-c)/b - Q_{-i} = 2m - Q_{-i} = 2x$. This is a concave parabola in $q_i$ with peak at $x$:
$\partial u_i/\partial q_i = b(2x - 2q_i)$, $\partial^2 u_i/\partial q_i^2 = -2b < 0$. With
$q_i \ge 0$,
\[
  BR_i(Q_{-i}) = \max\{0,\ x\} = \max\Big\{0,\ m - \frac{Q_{-i}}{2}\Big\}.
\]
Expanding, for any two outputs $a'$ and $t$,
\[
  u_i(a', Q_{-i}) - u_i(t, Q_{-i}) = b\big[2x(a' - t) - (a'^2 - t^2)\big]
  = b\,(t - a')(t + a' - 2x). \tag{1}
\]
So if every possible peak $x$ lies in $[\underline a, \bar a]$, outputs above $\bar a$ are
strictly dominated by $\bar a$ (both factors positive) and outputs below $\underline a$ by
$\underline a$ (both negative). An output that is the unique best reply to some surviving
$Q_{-i}$ is not dominated.

\textbf{(a) Two firms.} Here $Q_{-i} = q_j$. If $S_j^k = [L_k, U_k]$, the best replies to
$q_j \in [L_k, U_k]$ fill $[BR(U_k), BR(L_k)]$, so
\[
  L_{k+1} = \max\{0,\ m - \tfrac{U_k}{2}\}, \qquad U_{k+1} = m - \tfrac{L_k}{2}.
\]
\begin{itemize}
  \item Round~1: $q_j \in [0, \infty)$, so the peaks $x = m - q_j/2$ are all $\le m$ and the best
  replies fill $[0, m]$ ($BR = m$ at $q_j = 0$, $BR = 0$ for $q_j \ge 2m$). By (1) every
  $q_i > m$ is strictly dominated by $m$. $S_i^1 = [0, m]$.
  \item Round~2: $q_j \in [0, m]$ gives peaks in $[m - m/2, m] = [m/2, m]$; every $q_i < m/2$ is
  dominated by $m/2$. $S_i^2 = [m/2, m]$.
  \item Round~3: $q_j \in [m/2, m]$ gives $[m - m/2, m - m/4] = [m/2, 3m/4]$.
  Round~4: $[m - 3m/8, m - m/4] = [5m/8, 3m/4]$. Round~5: $[5m/8, 11m/16]$.
\end{itemize}
From round~2 on the bounds are interior ($U_k \le m < 2m$), so subtracting the recursions gives
$U_{k+1} - L_{k+1} = \tfrac12 (U_k - L_k)$: the width is $m/2^{k-1} \to 0$. Adding them, the
midpoint satisfies $M_{k+1} = m - M_k/2$, whose fixed point solves $q = m - q/2$, i.e.\
$q = 2m/3$, and $M_k - 2m/3 = (-\tfrac12)^{k-1}(M_1 - 2m/3) \to 0$. The nested intervals shrink to
one point:
\[
  S_i^\infty = \Big\{\tfrac{2m}{3}\Big\} = \Big\{\frac{a - c}{3b}\Big\}, \qquad i = 1, 2 .
\]
IESDS predicts the Cournot outcome $q_1 = q_2 = (a-c)/(3b)$, the slides' $\tfrac23 q^m$
(Lecture~2, slide~24; Tadelis's example has $m = 45$ and gives 30).

\textbf{(b) Three firms.} Now $Q_{-i} = q_j + q_h$.
\begin{itemize}
  \item Round~1: $Q_{-i} \in [0, \infty)$, so the best replies fill $[0, m]$ and every $q_i > m$
  is strictly dominated by $m$. $S_i^1 = [0, m]$.
  \item Round~2: now $Q_{-i} \in [0, 2m]$, so the peaks are $x = m - Q_{-i}/2 \in [0, m]$ and
  $BR(2m) = \max\{0, m - m\} = 0$. The best replies still fill all of $[0, m]$: each
  $q \in [0, m]$ is the unique best reply to $Q_{-i} = 2(m - q)$, which is reached by
  $q_j = q_h = m - q \in [0, m]$, two strategies that survived round~1. Nothing is deleted:
  $S_i^2 = S_i^1$.
\end{itemize}
Since round~2 deletes nothing, the process stops: $S_i^\infty = [0, (a-c)/(2b)]$ for each firm.
\emph{No}, IESDS does not give a unique prediction with three firms. The three-firm Cournot
equilibrium $q_i = (a-c)/(4b) = m/2$ (the $n$-firm exercise, Tadelis 5.8, with $n = 3$ and
$b = 1$) lies inside, but so does every other output in $[0, m]$.

\interp With two firms the best-reply map is a contraction (slope $-\tfrac12$): each round of ``I
know that you know \dots'' halves the range of plausible outputs, so common knowledge of
rationality alone pins down the Cournot outcome. With three firms the rivals' \emph{combined}
output can be as large as $2m$, enough to make zero output a best reply, so the first-round bounds
reproduce themselves and higher-order reasoning adds nothing. IESDS is weak in games with several
players whose actions are strategic substitutes; Nash equilibrium (Lecture~2) still pins down the
outcome, but only by assuming correct beliefs.

\verify \texttt{check\_A1.py} verifies with sympy that the Cournot payoff equals
$b q_i(2x - q_i)$ with $x = m - q_j/2$, the factorization (1), and that the fixed point of
$q = m - q/2$ is $(a-c)/(3b)$. It runs the duopoly interval recursion for 60 rounds in exact
arithmetic ($m = 1$: $[0,1]$, $[\tfrac12, 1]$, $[\tfrac12, \tfrac34]$, $[\tfrac58,\tfrac34]$,
$[\tfrac58,\tfrac{11}{16}]$, \dots, $\to \tfrac23$), checks that every $q \in [0, m]$ is the best
reply to $Q = 2(m - q)$ with both rivals in $[0, m]$, and runs brute-force IESDS for three firms on
a grid ($a - c = 12$, $b = 1$, steps of $0.5$), which leaves exactly $[0, 6] = [0, m]$.
\emph{Comparison with the MWG solutions manual} (pp.~8-3--8-4): it agrees on both parts,
$(a-c)/(3b)$ for two firms and $q_i \in [0, (a-c)/(2b)]$ with no unique prediction for three.

\mistakes
\begin{itemize}
  \item Stopping (a) at ``the best replies cross at $(a-c)/(3b)$''. That is the Nash argument; the
  exercise asks for dominance, so each round must name what is deleted and by what.
  \item In (b), bounding $Q_{-i}$ by $m$ instead of $2m$. Each rival can produce up to $m$, so the
  sum reaches $2m$, and that is what makes $q_i = 0$ a best reply.
  \item Concluding from (b) that the three-firm game has many Nash equilibria. It has one; only
  the dominance argument fails to find it.
\end{itemize}

\answer{(a) $BR_i(q_j) = \max\{0, (a-c)/(2b) - q_j/2\}$; IESDS gives $[0, m]$, $[m/2, m]$,
$[m/2, 3m/4]$, $[5m/8, 3m/4]$, \dots, with widths halving, so the unique survivor is
$q_1 = q_2 = (a-c)/(3b)$. (b) No: with three firms the rivals' total can reach $2m$, every
$q_i \in [0, (a-c)/(2b)]$ is a best reply to some surviving profile, and IESDS stops at
$[0, (a-c)/(2b)]$ for each firm.}
\end{exsol}
